%% P14 \dash{} Funkcje wielu zmiennych i gradient
%%
%% Zalążek. Poniższy opis jest kontraktem tego programu: co ma uzasadnić i co
%% ma dać czytelnikowi. Napisanie programu oznacza usunięcie bloku
%% \programstub{} i zastąpienie go programem. Jeśli po przerobieniu materiału
%% opis okaże się błędny, popraw go i odnotuj sprzeczność w CLAUDE.md w sekcji
%% *Resolved questions*.

\program{Funkcje wielu zmiennych i gradient}
\label{prog:P14}

\programstub{%
\textit{[Opis do przetłumaczenia \dash{} patrz wydanie angielskie.]}

Argument: a partial derivative holds everything else still; the gradient
collects them; the directional derivative is a dot product with the
gradient, so the gradient is the steepest direction and is perpendicular to
the level set. Payoff: gradient descent's direction is derived rather than
asserted; the zig-zag in a narrow valley is explained by the gradient being
perpendicular to the contour rather than pointing at the minimum, which is
the picture that makes momentum obvious two parts later.

\medskip
\textbf{Planowana długość:} 55 ramek.
}
